\subsection{The epistemic status}\label{sec:scale-status}

The program is either one proof away---H1g and the sharp H2m, both
reduced to concrete named objects, the strand census and the pair-line
count---or one systematic obstruction away: a family class whose
junction structure escapes the $\mathrm{ov}^{\,m-1}$ volume law. For
the obstruction direction the $T=8$ evidence points the other way: the
measured margin \emph{grows} with $k$, and no distinct family reaches
even a tenth of the chain-class volume at equal $\mathrm{ov}$ beyond
$k=2$. The distinction will be settled by whether the strand census
closes at $T=8$ with uniform constants, not by further measurement.

One more disclosure belongs in the record. The sequence of named
obligations has moved: the input structure, then the exhaustiveness of
the census classification, then the scaling-closure lemma itself, then
H1-general, now the strand census. Each closure has been real, and
each new obligation has been at the same scale as the previous one.
Whether this is a fractal obstruction---each gap revealing a new gap of
the same size, indefinitely---or a sequence of genuinely decreasing
obligations that will bottom out is not decidable from the present
data; after five iterations of the pattern the defensible prior is that
the reformulation has a stubborn residual, and Theorem~\ref{thm:SC} is
stated so that the residual is exactly what the hypotheses name. What
would change the verdict: a family at some $(T=8,k)$ whose per-size
volume exceeds $\mathrm{ov}^{\,m-1}$---the measured pool collapse
makes this increasingly unlikely, as admissibility itself dies at
$13\%$ by $p=29$---or a strand-census proof at $T=8$ with uniform
constants, which would make H1g a theorem at the rung and supply the
template for general $T$. No claim is made about the amount of work
needed to close the open obligations; the
failure modes are stated where the obligations are.

The record since that disclosure adds a sixth attack turn and a change
in the pattern's shape. The strand census has now been attacked twice
and has survived both attacks: the necessary-condition families of
this section (the overlap budget, the spacing caps, the last-arc fit
count---non-discriminating), and the marginal-rigidity route of
\S\ref{si:sec:scale-h1} (the projection envelope measured full, the two
naive coordinate readings false, the corrected character system closed
but silent on the count). Both closed as route-closures, and---for the
first time in the sequence's six attack turns---the named obligation
did not move: no successor lemma was minted; what moved was the
obligation's coordinate system (Lemma~\ref{si:lem:char10}) and its
numerical form (Proposition~\ref{si:prop:threshold}). The reading
has tightened accordingly: the residual is not a bookkeeping artifact
awaiting the right census format; it is resistant to the toolkit that
produced every proved result in this paper. A structural fact in the
ledger explains why that may be: the proved results are closures,
invariances, exhaustiveness statements, and machine verifications,
while every remaining obligation---H1g's margin over the trivial
bound, the volume law's unproved growth in $k$, the chain class's gap
from measured to proved constants---is a growth-in-$k$ margin over a
trivial or measured baseline, and no turn of this program has proved a
statement of that type. Both closed routes, moreover, were
census-shaped attacks on what Proposition~\ref{si:prop:threshold} shows
is an inequality-shaped obstruction. Three candidate third routes are
on record, none attempted here and none with a stated prior:
pairwise-difference marginals in the corrected $X$ coordinates; a
combinatorial characterization of the joint support, given its
measured sparsity ($1.5\%$--$7.4\%$ median density); and a
transfer-matrix or automaton model of the joint density under the
arc-covering constraint.

The type-mismatch reading of the preceding paragraph has since been
verified statement-by-statement against the paper's full proof
inventory---every lemma, proposition, theorem, and machine record---and
it holds, with the three statements that could be read as exceptions
resolved explicitly. First, the spacing lemma (Lemma~\ref{si:lem:spacing})
is a packing bound whose slack over its own trivial baseline is the
factor $b_{\min}$; at the cascade's operating point ($L=Rk$ against
$p=Tk$) the bulk-fiber value $b_{\min}\gtrsim p/(2L+1)\approx T/(2R)$ is
$k$-free, so the lemma is invariant under exactly the scaling in which
the obligations' margins grow, and the ``$k$-cancellation''
(Proposition~\ref{si:prop:kcanc}) is a $k$-invariance of the cut,
conditional on the volume law---the cancellation of the dependence the
obligations need to grow, not a proof that it grows. Second, part (d) of
the comparison (Proposition~\ref{si:prop:depth}) does state a margin
growing linearly in $k$, but as a consequence of its hypothesis stack
(H1g included), not as a proved statement about the measured world.
Third, the chain-class volume law (Lemma~\ref{si:lem:chainbox}) computes a
count that grows in $k$, but its margin over the trivial box is the
$k$-free threshold value of Proposition~\ref{si:prop:threshold}; the
growth-in-$k$ content of H1g is precisely the distinct class's extra
thinness below the chain benchmark ($10\times/43\times/160\times$ at
equal $\mathrm{ov}$, measured), which the chain-class theorem does not
touch, while the proved chain constants of Theorem~\ref{si:thm:chaindilate}
carry order factors ($4m\,s_{\max}$ with $s_{\max}\approx2k$) whose
slack grows in $k$---the measured-to-proved constant gap already listed
among the obligations. With the reading verified, the three candidate
third routes of the preceding paragraph are typed by output:
pairwise-difference marginals in the corrected $X$ coordinates and the
joint-support characterization are census-shaped (classification
outputs; the same type as the two routes already closed), while the
transfer-matrix or automaton model is the one margin-shaped
candidate---the natural output of a transfer-operator argument is a
spectral quantity uniform in $k$, the missing statement type. The typing
assigns output types, not success probabilities, and rests on the
tool-reach half of the reading rather than a theorem; it is the reason
the automaton route is the only one the record prioritizes, and the
attempt itself is its test.

The attempt has since been run, under the discipline of a proof
attempt: the two design specifications fixed beforehand, the success
criterion pre-committed, and the verdict taken against it. The route
closed by theorem \emph{in its single-anchor scope}
(Lemma~\ref{si:lem:monomial} and the analysis following it): every spectral
quantity the single-anchor automaton framework can produce is
exactly $0$ or $1$, the cyclic product's trace is the cascade's own
survivor count, and the survival-fraction deficit the criterion
demanded translates exactly into the pair bound H2m and the volume law
H1g---the instrument consumes the hypotheses it was commissioned to
prove. The typing's prediction was therefore right about the output
type and wrong about its shape: the route's natural output is not a
margin in disguise but a closure about itself, and the expected
spectral gap was an artifact of the classical analogy (sub-shifts of
finite type earn their gaps through branching that a lossless,
sufficient-statistic dynamics cannot have). Three routes into the
strand census are now closed---the necessary-condition families, the
marginal-rigidity route, and the (single-anchor) automaton---the first
two by measurement, the third by theorem \emph{within its stated
scope}. The empirical half of the reading (S1) is verified a second
time by a second instrument: the one margin-shaped candidate, fully
specified and executed, emitted exact census computations and a
$k$-invariance, and no margin. The tool-reach half (S2) has its first
proved instance---every automaton built on a bijective drift is
census-shaped---which is a mechanism, not yet the general claim;
instruments that branch (the general strand-census operator) or do not
inherit the drift remain outside its scope, and nothing here is claimed
about them. With this, the program's state is fixed, and it is the
state the conditional theorem was stated to name: Theorem~\ref{thm:SC}
with its ledger; the type-mismatch diagnosis, now with its mechanism
\emph{at the single-anchor level} (losslessness forbids generated decay
for that operator); and a named residual---the growth-in-$k$ margin,
uniform over families, size tuples, and $k$---that the single-anchor
toolkit is not the right shape to close, with the general strand-census
route left open. The reformulation closes the covering obstruction exactly
where classification suffices, and the conjecture-level consequence is
the one stated in the introduction: within this framework the residual
is a structural gap for the single-anchor toolkit, not a technical one
for that toolkit; closing it in general is work for a different
instrument, and the general strand-census route is left open.

The disposition of that verdict is now on record, and it fixes the
program's final state. The natural external instrument for the missing
margin---the Ehrhart theory of the Lonely-Runner zonotope, the third
bridge of \S\ref{sec:open}---was engaged as a full strand and executed
to certificates. Its arithmetic is unimpeachable: the $h^*$-polynomials
of both canonical families are real-rooted, log-concave, and Newton
through $n=10$, in exact rational arithmetic throughout. Its geometric
reading is false: the deepest-hole form fails from $n=5$ with exact
witnesses (the harmonic instance closed \emph{exactly} in the fourth
revision, $\rho(1,2,3,4,5)=16/47$; the two $n=6$ boundary rows
resolved as strict refutations, $\rho\ge17/47$ at $m=12$ and
$\rho\ge1793/5000$ at $m=9$),
and the modular law it suggested---deepest hole at the center exactly
when $n\mid m$---holds exactly at $n=3$ ($m\le 24$), is certified at
$n=4$ (through $m=16$, value $9/34$), breaks at $n=5$ (the harmonic
$m=5$ and $m=15$ refuted exactly, $\rho\ge 97/288$ and $6/19$; the
surviving multiples $m=10,20,25,30$ certified at $7/22$, $13/42$,
$4/13$, and $19/62$ respectively---$m=20,25,30$ in the first
certification and $m=10$ closed by the companion paper's
third attempt, which paired a heavily scaled random-local
search with an exact Lipschitz box certificate and emptied the tree),
and is
extinct at $n=6$ (every tested instance refuted exactly, through the
$4n$ dilate, the narrowest margin $71/208$ against $17/50$). The
instrument that most faithfully preserves the instance's arithmetic is
again the one that cannot produce the margin: the type mismatch is a
property of the lossless encodings themselves, not of one toolkit.
The record is accordingly published in two parts---this monograph (the
exact identities, the classifications, the conditional theorem with
its ledger) and a companion paper \cite{companion-paper} that works
out the negative half: the four falsified strengthenings, the witness
archive, the certification ladder above with its negative control, and
the $h^*$ addendum. The scope statement of the abstract applies to
both parts unchanged.
